\section{Data and measures}\label{sec:setup}

Which flights, which populations, and which measures of the motion?

\paragraph{Flights and groups.} 156\,391 paraglider and 6\,094 hang-glider flights
 on 6\,209 days, the cleaned tracks of
the old repository (soaring, tag \texttt{legacy-final}). A flight is analysed after its
\emph{departure}, the first fix 3~km from the takeoff; before it the pilot counts as neither solo
nor in a group. The measures of the motion use loggers of at most 10~s; the phases (climb, search,
glide) are Vilpellet's runs and exist on the flights logged every second, 70\% of the paragliders.
The groups are those of the note, with $\delta=2$~min, $\varepsilon_h=500$~m, $\varepsilon_z=300$~m,
$\theta_{\max}=45^\circ$, $g=2$~min and $q=0.1$, hence contacts of at least $\tau^*=\tauStar$~min
(its sections 3--6). A contact means shared air within sight; it does not prove interaction.

\paragraph{Populations.} Three populations of stretches of flight.
\begin{itemize}
  \item $\Sl$, \emph{solo flights}: flights never in a group.
  \item $\G$, \emph{group time}: the stretches of a flight inside a configuration, a fixed set of
        flights that stays one group over a maximal interval. $\G_n$ is the group time in
        configurations of $n$ pilots; ``$n$ pilots'' always means the size of the configuration at
        that instant.
  \item $\A$, \emph{the same flights alone}: the stretches of the group flights outside any
        configuration. It has the pilots, gliders, loggers and days of $\G$, and serves as its
        control.
\end{itemize}

\paragraph{The motion over every lag.} The horizontal position $\mathbf r(t)$ of a flight is read on
a 10-s grid of its clock. For a flight $f$ and a population, the time-averaged mean squared
displacement pools the origins of all its stretches in that population,
\begin{equation}\label{eq:tamsd}
  \overline{\delta^2_f}(\tau)=\frac{1}{N_f(\tau)}\sum_{t}\bigl\|\mathbf r(t+\tau)-\mathbf r(t)\bigr\|^2 ,
\end{equation}
where $t$ runs over the $N_f(\tau)$ origins with $t$ and $t+\tau$ in the same stretch, and the
ensemble $\MSD(\tau)$ is the mean of $\overline{\delta^2_f}(\tau)$ over the flights, one weight per
flight. A stretch enters only if it lasts at least 1250~s, so that every population keeps the same
flights at all 21 lags from 10 to 1000~s. With $\MSD(\tau)\propto\tau^{2H}$, $H$ is half the slope
of $\log\MSD$ against $\log\tau$ over all lags: $H=1$ for a straight line, $H=1/2$ for a random walk.
$\Hloc(\tau)$ is the same slope over the lags within $\pm0.25$ decades of $\tau$
(Appendix~\ref{app:hloc}). We read it at two scales: $\tau=100$~s, within one climb or one glide,
and $\tau=1000$~s, over several climb--glide cycles.

\paragraph{Climbs and glides.} Each phase run goes whole to the population of the stretch that holds
its middle, and each statistic is a ratio of sums over the runs of a population, so long phases weigh
more. Of the twelve statistics measured, five are not explained by the others and are the ones drawn
(Appendix~\ref{app:redundancy}): the share of time in climb, the share of searches that end in a
climb, the climb rate (total height gained over total climb time), the glide speed and the glide
ratio, replaced in stratified contrasts by its inverse, the height lost per kilometre of glide
(Appendix~\ref{app:hloc}). The tables also give the cross-country speed, the straight-line distance
covered in 30~min averaged over every origin, which the five explain.

\paragraph{Strata.} A \emph{zone} is a cell of 25~km containing the middle of a piece of track; a
\emph{site-day} holds the flights of one day whose takeoffs lie within 3~km; the \emph{hour} is the
local hour of the middle of a piece or a run.

\paragraph{Uncertainty.} Bands are 5--95\% of bootstrap draws that resample whole dependency
blocks, the flights tied by sharing episodes, and recompute the statistic: the blocks of the note
for pooled statistics (1000 draws), blocks that also join the flights of a site-day for stratified
ones (200 draws, the same for every statistic). They measure the sampling of blocks, not the choice
of $q$ or of the tolerances. Two bands that do not overlap are not a test; differences are computed
draw by draw.

\paragraph{What the grouping gives.} Of the 154\,598 paraglider flights with a logger of at most 10~s, 33\,501 (21.7\%) fly in a group at some point; of the 6\,039 hang-glider flights, 418 (6.9\%). Group time is 12.8\% of all paraglider flight time after the departure, and on average 45\% of the time of a group flight. The 36\,575 pair intervals last 40 to 96~min (10th to 90th percentile; median 52~min) and form 40\,051 groups in 17\,169 episodes. Half of the episodes have two flights, 1\% have more than 9, the largest has 53. By the size of the group, the paraglider group time is: 2 pilots 64\%, 3 pilots 19\%, 4 pilots 8\%, 5-6 pilots 6\%, 7+ pilots 3\% (the largest group has 27 pilots). In contact, one pilot is behind the other by a median of 20~s (75\% within 45~s).


\paragraph{What we keep.} Three populations, $\Sl$, $\G$ and $\A$; two scales of the motion,
$\Hloc(100~\mathrm{s})$ and $\Hloc(1000~\mathrm{s})$; the statistics of climbs and glides; bands
from blocks of flights that are not independent.
